\documentclass[conference]{IEEEtran}

% ============================================================
% PACKAGES
% ============================================================
\usepackage[utf8]{inputenc}
\usepackage[T1]{fontenc}
\usepackage{cite}
\usepackage{amsmath,amssymb,amsfonts}
\usepackage{graphicx}
\usepackage{textcomp}
\usepackage{xcolor}
\usepackage{booktabs}
\usepackage{array}
\usepackage{multirow}
\usepackage{url}
\usepackage{hyperref}
\usepackage{float}
\usepackage{enumitem}
\usepackage{balance}
\usepackage{tabularx}

\hypersetup{
    colorlinks=true,
    linkcolor=black,
    citecolor=black,
    urlcolor=blue
}

\setlist{nosep, leftmargin=*}

\begin{document}

% ============================================================
% TITLE
% ============================================================
\title{VisionGuard AI: A Lightweight, Explainable Multi-Task Deep Learning Pipeline for Automated Glaucoma Screening from Retinal Fundus Images}

\author{
    \IEEEauthorblockN{Muhammad Ijlal Asif\textsuperscript{1} (F24607072), \quad Syed Maaz Shah\textsuperscript{2} (F24607091)}
    \IEEEauthorblockA{
        Department of Artificial Intelligence (2024)-B \\
        National University of Technology (NUTECH), Islamabad, Pakistan \\
        Subject: Machine Learning (CS-284) \quad|\quad Submitted To: Ma'am Faiza
    }
}

\maketitle

% ============================================================
% ABSTRACT
% ============================================================
\begin{abstract}
Glaucoma is a progressive optic neuropathy and one of the leading causes of irreversible blindness worldwide, affecting over 80 million individuals globally. Early detection through retinal fundus image analysis is critical for preventing permanent vision loss, yet manual screening by trained ophthalmologists remains slow, expensive, and largely inaccessible in resource-limited regions. This paper presents VisionGuard AI, a complete research-to-deployment deep learning pipeline for automated Glaucoma screening that was developed iteratively across four project phases. In Phase~1, a baseline four-layer Convolutional Neural Network achieved only 72.5\% accuracy due to limited feature extraction capacity. Phase~2 introduced Transfer Learning using EfficientNet-B0 pre-trained on ImageNet, raising accuracy to 99.05\% on the ACRIMA validation set. Phase~3 integrated Grad-CAM explainability to address the clinical ``black-box'' problem and delivered a Streamlit web interface for interactive use. In this final Phase~4, the architecture was fundamentally restructured into a Teacher-Student Knowledge Distillation framework. The Teacher model, a ResNet-50 backbone enhanced with a Convolutional Block Attention Module~(CBAM), is trained using a Joint Loss function that combines Cross-Entropy Loss for classification with Dice Loss and Binary Cross-Entropy Loss for simultaneous optic cup and disc segmentation. The Student model, a lightweight EfficientNet-B0 with CBAM, is subsequently distilled from the Teacher using a composite loss of ground-truth supervision, KL-Divergence soft-label transfer at temperature $T{=}4.0$, and segmentation mask mimicry. The production system incorporates Test-Time Augmentation~(TTA) for robust predictions, Out-of-Distribution~(OOD) detection via normalized Shannon Entropy, dual Explainable AI visualizations through Grad-CAM and Integrated Gradients, a Human-in-the-Loop adjustable diagnostic threshold, and Temperature Scaling~($T{=}2.5$) for probability calibration. The final system is deployed as a FastAPI REST backend consumed by a React-based clinical dashboard, achieving a deployed model size of 24\,MB with sub-200\,ms inference time on CPU hardware, making it practically suitable for real-world clinical screening environments in underserved regions.
\end{abstract}

\begin{IEEEkeywords}
Glaucoma Detection, Knowledge Distillation, CBAM, Multi-Task Learning, Explainable AI, Grad-CAM, Integrated Gradients, Transfer Learning, Fundus Images, Medical Imaging.
\end{IEEEkeywords}

% ============================================================
% I. INTRODUCTION
% ============================================================
\section{Introduction and Final Problem Statement}

\subsection{Problem Definition}
Glaucoma is a group of progressive eye diseases that damage the optic nerve head through elevated intraocular pressure, leading to permanent and irreversible vision loss. According to the World Health Organization, over 80 million people suffer from Glaucoma globally, with projections reaching 112 million by 2040~\cite{who2019}. The condition is frequently termed ``the silent thief of sight'' because it progresses entirely asymptomatically in its early stages; a patient can lose up to 40\% of their peripheral visual field before experiencing any noticeable symptoms. By the time symptoms become apparent, significant and irreversible damage to the retinal nerve fiber layer has already occurred. The primary clinical indicator of Glaucoma is an elevated vertical Cup-to-Disc Ratio~(CDR), where the optic cup --- the central depression within the optic disc --- becomes abnormally enlarged relative to the surrounding disc boundary due to progressive death of retinal ganglion cell axons.

\subsection{Real-World Importance}
Currently, diagnosing Glaucoma requires a trained ophthalmologist to manually examine retinal fundus photographs using specialized equipment, visually estimate the CDR, and identify subtle signs such as neuroretinal rim thinning, disc hemorrhages, and nerve fiber layer defects. This diagnostic process is inherently slow, subjective, and heavily dependent on the individual clinician's experience and training. Studies have shown significant inter-observer variability even among expert ophthalmologists when estimating CDR values. In developing countries such as Pakistan, the ratio of eye care specialists to the general population is severely inadequate, particularly in rural and semi-urban areas where access to ophthalmic care may require traveling hundreds of kilometers. As a direct consequence, millions of patients are never screened for Glaucoma and lose their vision permanently without ever receiving a diagnosis. An automated, accurate, and --- critically --- explainable AI screening tool that can operate on standard consumer-grade hardware without requiring specialized medical imaging infrastructure would directly address this life-altering healthcare gap.

\subsection{Final Scope of the Project}
VisionGuard AI delivers a fully trained, lightweight, and clinically transparent deep learning system accessible through a modern web-based clinical dashboard. The system performs multi-task analysis on uploaded retinal fundus images, simultaneously producing a binary classification (Normal vs.\ Glaucomatous) and generating pixel-level optic cup and disc segmentation masks for CDR computation. Beyond raw prediction, the system provides dual Explainable AI~(XAI) visualizations through two complementary methods: Grad-CAM for coarse regional localization of diagnostically relevant areas, and Integrated Gradients for fine-grained pixel-level attribution maps. These visualizations allow clinicians to visually verify and trust the AI's reasoning before acting on its recommendation. Additional clinical safety features include Out-of-Distribution~(OOD) image rejection to prevent unreliable predictions on non-fundus inputs, an adjustable diagnostic threshold slider enabling Human-in-the-Loop calibration of the sensitivity-specificity tradeoff, and Test-Time Augmentation~(TTA) to stabilize predictions against minor variations in image orientation and capture conditions.

\subsection{Summary of the System Evolution}
The project evolved through four iterative development phases, each addressing specific limitations identified in the previous phase. Phase~1 established a performance baseline using a simple custom CNN architecture. Phase~2 introduced Transfer Learning with modern ImageNet-pretrained architectures, dramatically improving accuracy. Phase~3 addressed the ``black-box'' problem that is the primary barrier to clinical adoption by integrating visual explainability and delivering a web-based interface for interactive demonstrations. Phase~4, fully documented in this report, represents the final and most significant transformation: restructuring the entire system into a Teacher-Student Knowledge Distillation framework with bottleneck attention mechanisms, multi-task learning with a U-Net-style segmentation decoder, dual XAI outputs, clinical safety mechanisms, and a production-grade React frontend communicating with a FastAPI REST backend.

% ============================================================
% II. LITERATURE REVIEW
% ============================================================
\section{Literature Review}

Automated Glaucoma detection using deep learning has emerged as one of the most rapidly growing research areas in ophthalmic artificial intelligence. A recent bibliometric scoping review confirmed that over 11\% of all historical publications in this field were produced in the last four years alone, underscoring the urgency and relevance of this research direction. This section reviews the key studies that directly influenced our architectural decisions and positions VisionGuard relative to the existing body of published work.

Li~et~al.~\cite{li2018} developed one of the earliest large-scale deep learning systems for Glaucoma screening using color fundus photographs. Their ResNet-based architecture, trained on a proprietary dataset of over 48,000 images from multiple clinical sites, achieved an impressive Area Under the Curve~(AUC) of 0.986 and demonstrated that automated AI performance can match or exceed that of experienced, board-certified ophthalmologists. While their results were groundbreaking, their dataset remained proprietary and unavailable to the research community, severely limiting reproducibility and independent validation. Our work specifically addresses this gap by utilizing the publicly accessible ACRIMA dataset, ensuring that every aspect of our pipeline can be independently reproduced and verified.

Diaz-Pinto~et~al.~\cite{diazpinto2019} created the ACRIMA dataset, which has since become a widely used benchmark in the Glaucoma detection community. The dataset contains 705 expert-labeled retinal fundus images annotated by experienced Glaucoma specialists. In their original study, they benchmarked several classical CNN architectures including VGG-16, ResNet-50, and GoogLeNet, reporting classification accuracies exceeding 95\%. However, their evaluation was limited exclusively to single-task classification without any explainability component, segmentation output, or deployment consideration. Our project extends their foundational work by applying a modern Teacher-Student distillation framework with CBAM attention and dual XAI outputs on the same dataset, significantly advancing the state of clinical readiness.

Orlando~et~al.~\cite{orlando2020} organized the REFUGE Challenge, which established the first large-scale standardized benchmark for comprehensive Glaucoma assessment from fundus photographs, including both classification and optic disc/cup segmentation tasks. Their systematic analysis of submitted methods revealed a critical insight: ensemble methods and multi-task approaches that jointly perform segmentation alongside classification significantly and consistently outperform single-task classifiers. This key finding from the challenge results directly motivated our decision to adopt a Multi-Task Learning architecture that simultaneously performs classification and optic disc/cup segmentation, leveraging the synergy between these two clinically related tasks.

Tan and Le~\cite{tan2019} introduced the EfficientNet family of convolutional neural networks, which achieve superior accuracy-to-parameter ratios through a novel compound scaling method that uniformly scales network depth, width, and input resolution using a fixed set of scaling coefficients. EfficientNet-B0, the smallest variant in the family with only 5.3 million parameters, outperforms considerably larger older architectures such as VGG-16 (138M parameters) and even ResNet-50 (25.6M parameters) on the ImageNet benchmark while consuming a fraction of the computational budget. This exceptional efficiency-accuracy tradeoff made EfficientNet-B0 the natural and ideal candidate for our deployed Student model, which must run efficiently on resource-constrained hardware.

He~et~al.~\cite{he2016} introduced the Residual Network~(ResNet) architecture with its signature skip connections that enable effective training of very deep networks by mitigating the gradient vanishing and degradation problems. ResNet-50 has become one of the most widely adopted backbone architectures in medical imaging research due to its strong, generalizable feature extraction capabilities and proven track record across diverse medical imaging modalities. We leverage ResNet-50 as our Teacher model specifically because its deeper architecture and higher parameter count enable it to capture richer, more nuanced feature representations of the optic nerve head, which can subsequently be effectively distilled into the lighter Student model.

Selvaraju~et~al.~\cite{selvaraju2017} proposed Gradient-weighted Class Activation Mapping~(Grad-CAM), a post-hoc visualization technique that produces visual explanations from deep networks by computing the gradient of the target class score with respect to feature maps of the final convolutional layer. The resulting weighted combination of feature maps is upsampled to input resolution, producing a coarse localization heatmap highlighting the regions most influential to the model's prediction. In clinical AI systems, explainability is not merely a desirable feature --- it is increasingly mandated by regulatory frameworks and demanded by medical professionals who must bear legal responsibility for diagnostic decisions. We integrate Grad-CAM as our primary XAI method, specifically targeting the CBAM attention module to generate clinically meaningful, anatomically-grounded heatmaps.

Sundararajan~et~al.~\cite{sundararajan2017} introduced Integrated Gradients, an axiomatic attribution method that satisfies two fundamental theoretical properties: \textit{Sensitivity} (if changing a single feature changes the prediction, that feature must receive a non-zero attribution) and \textit{Implementation Invariance} (attributions are guaranteed to be identical for any two functionally equivalent networks, regardless of their internal implementation). Unlike Grad-CAM, which provides coarse, region-level explanations, Integrated Gradients operates at the individual pixel level by numerically integrating the model's gradients along a straight-line path in input space from a neutral baseline image (typically a black image) to the actual input. We implement this as our secondary, complementary XAI method using Meta's Captum library with 50 interpolation steps, providing fine-grained pixel-level attributions that give clinicians a second, independent perspective for verifying model reasoning.

Woo~et~al.~\cite{woo2018} proposed the Convolutional Block Attention Module~(CBAM), a lightweight and general-purpose attention mechanism that can be seamlessly inserted into any CNN architecture to improve representational power with minimal computational overhead. CBAM sequentially applies two complementary sub-modules: channel attention, which models inter-channel relationships to identify \textit{what} features are most discriminative, and spatial attention, which models inter-spatial relationships to identify \textit{where} the most informative features are located within the feature map. We integrate CBAM at the bottleneck of both our Teacher and Student architectures to force the networks to concentrate their representational capacity on the clinically relevant optic nerve head region, suppressing irrelevant background features.

Hinton~et~al.~\cite{hinton2015} formalized the concept of Knowledge Distillation, a model compression technique where a smaller, more efficient ``Student'' network is trained to replicate the behavior of a larger, more accurate ``Teacher'' network. The key insight is that the Teacher's softened probability distribution, produced using a temperature-scaled softmax, contains rich ``dark knowledge'' about inter-class similarities and relationships that is not captured by hard one-hot labels. By training the Student to match these soft distributions, knowledge is effectively transferred from the heavy Teacher to the lightweight Student. Our distillation pipeline extends this foundational concept to multi-task learning by simultaneously distilling both the Teacher's classification logits and segmentation masks into the Student model.

\smallskip
\noindent\textbf{Positioning of VisionGuard:} The vast majority of published Glaucoma detection systems rely on single-task classification using a single CNN backbone with, at most, one XAI visualization method. VisionGuard AI distinguishes itself by uniquely integrating six orthogonal technical contributions into a single, cohesive, end-to-end pipeline: (1)~Teacher-Student Knowledge Distillation for model compression, (2)~Multi-Task Learning with joint classification and segmentation, (3)~CBAM attention at the network bottleneck, (4)~dual XAI visualizations combining Grad-CAM with Integrated Gradients, (5)~clinical safety features including OOD detection, adjustable thresholds, and TTA, and (6)~full-stack production deployment via React and FastAPI. To our knowledge, no existing published work combines all six of these elements for Glaucoma screening.

% ============================================================
% III. PROJECT EVOLUTION
% ============================================================
\section{Complete Project Evolution}

\subsection{Phase 1: Baseline CNN (Assignment 1)}
The initial approach involved designing and training a custom four-layer Convolutional Neural Network entirely from scratch on the ACRIMA dataset. The architecture consisted of four convolutional blocks, each containing a convolution layer, batch normalization, ReLU activation, and max pooling, followed by fully connected layers for binary classification. This baseline model achieved only 72.5\% classification accuracy. The poor performance was attributed to three primary factors: (1)~the network's limited depth and parameter count provided insufficient feature extraction capacity for the nuanced visual patterns associated with Glaucoma, (2)~the absence of any data augmentation strategy led to severe overfitting on the small 705-image dataset, and (3)~training from randomly initialized weights, without any pre-training, meant the network had to learn all visual representations from scratch using extremely limited data.

\subsection{Phase 2: Transfer Learning (Assignment 2)}
To address the fundamental feature extraction limitations identified in Phase~1, we replaced the custom CNN with EfficientNet-B0 pre-trained on the ImageNet dataset containing 1.2 million natural images across 1,000 categories. The pre-trained backbone layers were initially frozen, and only the final classification head was replaced with a custom two-class linear layer and retrained on our Glaucoma data. Training used the Adam optimizer with a learning rate of 0.001 and Cross-Entropy Loss for 15 epochs with a batch size of 32. This Transfer Learning approach dramatically improved performance to 99.05\% validation accuracy. Data augmentation was introduced in the form of random horizontal flips and random rotation of up to $\pm$10\textdegree\ to combat overfitting. However, the system still functioned as a pure black-box binary classifier --- it produced only a ``Normal'' or ``Glaucoma'' label with no visual evidence, no test set evaluation had been performed, and no deployment interface existed.

\subsection{Phase 3: Explainability and Web Interface (Assignment 3)}
Phase~3 systematically addressed three critical gaps identified in Phase~2. First, we completed a comprehensive evaluation on the held-out test set of 107 images that the model had never seen during training or validation, confirming 99.05\% test accuracy with an AUC-ROC of 1.0 and an F1-Score of 0.99. Second, we integrated the pytorch-grad-cam library to generate Grad-CAM heatmap visualizations, which visually confirmed that the model consistently and correctly focuses its attention on the optic disc region --- the anatomical area that is clinically relevant for Glaucoma assessment. Third, we built a functional Streamlit web application that allows users to upload individual fundus images and receive real-time classification results alongside Grad-CAM heatmap overlays. While Phase~3 significantly improved the system's clinical transparency and usability, the underlying model remained a single-task classifier that was too heavy for mobile deployment and lacked robustness features such as OOD detection or prediction stabilization.

\subsection{Phase 4: Research-to-Deployment (Final --- This Report)}
The final phase represents the most significant architectural transformation of the entire project lifecycle. Every major component was either replaced or substantially enhanced:

\begin{itemize}
    \item \textbf{Data Preprocessing:} Enhanced with ColorJitter augmentation (brightness $\pm$0.2, contrast $\pm$0.2), random vertical flips, and increased rotation range ($\pm$15\textdegree) to better simulate real-world capture variability.
    \item \textbf{Model Architecture:} Completely restructured from a single EfficientNet-B0 classifier to a dual-model Teacher-Student Knowledge Distillation framework. Both Teacher (ResNet-50) and Student (EfficientNet-B0) were augmented with CBAM attention modules and U-Net-style segmentation decoders with skip connections.
    \item \textbf{Feature Engineering:} Shifted from single-task binary classification to Multi-Task Learning. The shared convolutional backbone simultaneously serves a classification head and a segmentation decoder, forcing the network to learn feature representations that encode both diagnostic class boundaries and spatial optic disc/cup structure.
    \item \textbf{Training Strategy:} Evolved from simple fine-tuning to a two-stage training process: (Stage~1) train the heavy Teacher model with a Joint Loss function, then (Stage~2) freeze the Teacher and distill its knowledge into the lightweight Student using KL-Divergence soft-label transfer and segmentation mask mimicry.
    \item \textbf{Inference Pipeline:} Added TTA (3-variant averaging), OOD detection (entropy thresholding), Temperature Scaling (calibrated probabilities), and a Human-in-the-Loop adjustable diagnostic threshold.
    \item \textbf{Deployment:} Migrated from a simple Streamlit prototype to a professional-grade FastAPI REST backend with a React-based clinical dashboard featuring glassmorphism UI design, interactive diagnostic controls, and real-time XAI visualizations.
\end{itemize}

% ============================================================
% IV. FINAL METHODOLOGY
% ============================================================
\section{Final Methodology (Optimized Pipeline)}

\subsection{Data Preprocessing and Feature Engineering}
All input fundus images are resized to $224 \times 224$ pixels to match the standard input dimensions expected by both backbone networks. Pixel values are converted to floating-point tensors and normalized using ImageNet channel statistics (mean $\mu = [0.485, 0.456, 0.406]$, standard deviation $\sigma = [0.229, 0.224, 0.225]$). During training, the following augmentation pipeline is applied to artificially expand the effective dataset size, introduce controlled variability, and prevent overfitting:

\begin{enumerate}
    \item \textbf{Random Horizontal Flip} (50\% probability) --- simulates left-right capture orientation.
    \item \textbf{Random Vertical Flip} (50\% probability) --- simulates vertical orientation variations.
    \item \textbf{Random Rotation} ($\pm$15\textdegree) --- simulates camera angle misalignment during fundus capture.
    \item \textbf{Color Jitter} (brightness $\pm$0.2, contrast $\pm$0.2) --- simulates varying illumination conditions across different fundus cameras and clinical environments.
\end{enumerate}

These transforms collectively simulate the natural variations encountered in real-world clinical image acquisition, including differences in camera orientation, lighting conditions, patient positioning, and imaging hardware. During validation and inference, only deterministic resizing, tensor conversion, and normalization are applied to ensure reproducible evaluation.

\subsection{Final Model Architecture}
The final architecture employs a Teacher-Student Knowledge Distillation framework where both models share a common architectural pattern: a pre-trained backbone, a CBAM attention module at the bottleneck, a classification head, and a U-Net-style segmentation decoder.

\subsubsection{Teacher Model (ResNet-50 + CBAM)}
The Teacher utilizes a ResNet-50 backbone pre-trained on ImageNet (using \texttt{ResNet50\_Weights.DEFAULT}). The backbone's hierarchical feature extraction produces multi-scale representations at four resolution levels: \texttt{layer1} (256 channels, $56{\times}56$), \texttt{layer2} (512 channels, $28{\times}28$), \texttt{layer3} (1024 channels, $14{\times}14$), and \texttt{layer4} (2048 channels, $7{\times}7$). A full CBAM module with channel reduction ratio $r{=}16$ and spatial convolution kernel size $k{=}7$ is applied to the 2048-channel bottleneck features, sequentially refining them through channel and spatial attention. The refined features feed into two task-specific heads:

\noindent\textit{Classification Head:} AdaptiveAvgPool2d$(1,1)$ $\rightarrow$ Flatten $\rightarrow$ Linear(2048, 2), producing two class logits for Normal and Glaucoma.

\noindent\textit{Segmentation Decoder:} A four-stage U-Net-style decoder with transposed convolutions and skip connections from encoder layers. Each decoder stage upsamples features by $2{\times}$ using \texttt{ConvTranspose2d}, concatenates with the corresponding encoder skip features along the channel dimension, and applies a $3{\times}3$ \texttt{Conv2d} with Batch Normalization and ReLU activation. The first stage takes the 2048-channel bottleneck output, upsamples to $14{\times}14$, and concatenates with \texttt{layer3} skip features (1024 channels), producing 1536 channels that are convolved down to 512. Subsequent stages progressively reduce channel width (512$\rightarrow$256$\rightarrow$128$\rightarrow$32) while increasing spatial resolution. The final stage uses a $4{\times}$ transposed convolution to reach full $224{\times}224$ resolution, followed by a $1{\times}1$ convolution producing a 2-channel output: channel~0 = optic disc probability mask, channel~1 = optic cup probability mask.

\subsubsection{Student Model (EfficientNet-B0 + CBAM)}
The Student follows the same multi-task dual-head pattern but uses the more efficient EfficientNet-B0 as its backbone (using \texttt{EfficientNet\_B0\_Weights.DEFAULT}). The backbone consists of 9 mobile inverted bottleneck~(MBConv) stages. Skip features for the decoder are extracted at stage~2 (24 channels, $H/4$), stage~3 (40 channels, $H/8$), and stage~5 (112 channels, $H/16$), with the bottleneck at stage~8 (1280 channels, $H/32$). CBAM ($r{=}16$, $k{=}7$) is applied at the 1280-channel bottleneck. The classification head maps Linear(1280, 2). The segmentation decoder uses proportionally smaller channel widths (1280$\rightarrow$256$\rightarrow$128$\rightarrow$64$\rightarrow$32$\rightarrow$16) with the same skip-connection architecture, producing a $224{\times}224{\times}2$ output.

\subsubsection{CBAM Attention Module}
CBAM~\cite{woo2018} consists of two sequential sub-modules that refine feature maps along complementary dimensions:

\noindent\textit{Channel Attention:} Global Average Pooling and Global Max Pooling are applied in parallel to the input feature map $\mathbf{F} \in \mathbb{R}^{C \times H \times W}$, each producing a $C{\times}1{\times}1$ descriptor. Both descriptors are independently passed through a shared two-layer MLP implemented as $1{\times}1$ convolutions (reducing channels from $C$ to $C/r$ and back, where $r{=}16$), and the outputs are summed and activated with Sigmoid:
\begin{equation}
    \mathbf{M}_c(\mathbf{F}) = \sigma\big(\text{MLP}(\text{AvgPool}(\mathbf{F})) + \text{MLP}(\text{MaxPool}(\mathbf{F}))\big)
\end{equation}

\noindent\textit{Spatial Attention:} After channel refinement ($\mathbf{F}' = \mathbf{M}_c \odot \mathbf{F}$), channel-wise average and max values are computed and concatenated into a 2-channel feature, which is convolved with a $7{\times}7$ kernel and activated with Sigmoid:
\begin{equation}
    \mathbf{M}_s(\mathbf{F}') = \sigma\big(\text{Conv}^{7\times7}([\text{AvgPool}_c(\mathbf{F}');\, \text{MaxPool}_c(\mathbf{F}')])\big)
\end{equation}

The final refined output is: $\mathbf{F}'' = \mathbf{M}_s(\mathbf{F}') \odot \mathbf{F}'$.

\subsection{Training Strategy and Loss Functions}

Training proceeds in two stages on a Google Colab NVIDIA T4 GPU with 16\,GB VRAM.

\subsubsection{Stage 1 --- Teacher Training}
The Teacher is trained for 15 epochs using Adam optimizer (learning rate $= 1{\times}10^{-4}$, weight decay $= 1{\times}10^{-5}$, batch size $= 16$). The Joint Loss function combines classification and segmentation objectives:
\begin{equation}
    \mathcal{L}_{\text{teacher}} = \alpha \cdot \mathcal{L}_{\text{CE}} + \beta \cdot (\mathcal{L}_{\text{Dice}} + \mathcal{L}_{\text{BCE}})
\end{equation}
where $\alpha = 1.0$ and $\beta = 1.0$. The Cross-Entropy Loss handles the binary classification task, while the Dice Loss (with smooth factor $\epsilon = 10^{-6}$) and Binary Cross-Entropy with Logits Loss jointly supervise the segmentation output. The Dice Loss is defined as:
\begin{equation}
    \mathcal{L}_{\text{Dice}} = 1 - \frac{1}{C}\sum_{c=1}^{C}\frac{2\sum_{i}p_{c,i} \cdot g_{c,i} + \epsilon}{\sum_{i}p_{c,i} + \sum_{i}g_{c,i} + \epsilon}
\end{equation}
where $p_{c,i}$ and $g_{c,i}$ represent the predicted and ground-truth values for class $c$ at pixel $i$.

\subsubsection{Stage 2 --- Student Training via Knowledge Distillation}
The frozen Teacher generates soft probability targets for each training batch. The Student is trained for 15 epochs using Adam (learning rate $= 2{\times}10^{-4}$, weight decay $= 1{\times}10^{-5}$, batch size $= 16$). The total distillation loss combines three components:
\begin{equation}
    \mathcal{L}_{\text{total}} = \mathcal{L}_{\text{GT}} + 0.5 \cdot \mathcal{L}_{\text{KD}}^{\text{cls}} + 0.5 \cdot \mathcal{L}_{\text{KD}}^{\text{seg}}
\end{equation}

The ground-truth loss $\mathcal{L}_{\text{GT}}$ uses the same Joint Loss structure but with $\alpha{=}0.5$ and $\beta{=}1.0$, reducing classification weight to emphasize segmentation learning during distillation. The classification KD loss is KL-Divergence between temperature-scaled softmax distributions:
\begin{equation}
    \mathcal{L}_{\text{KD}}^{\text{cls}} = T^2 \cdot \text{KL}\Big(\log\sigma\!\left(\frac{\mathbf{z}_s}{T}\right) \Big\| \sigma\!\left(\frac{\mathbf{z}_t}{T}\right)\Big)
\end{equation}
where $\mathbf{z}_s$ and $\mathbf{z}_t$ are Student and Teacher logits, and the distillation temperature $T{=}4.0$. The multiplication by $T^2$ compensates for the reduced gradient magnitudes caused by softening. The segmentation KD loss is the Dice Loss between the Student's segmentation output and the sigmoid-activated Teacher segmentation output, encouraging the Student to replicate the Teacher's spatial understanding of optic disc and cup boundaries. The checkpoint with the lowest validation loss is saved as the final model weights.

\subsection{Inference-Time Features}

\textbf{Test-Time Augmentation (TTA):} During inference, the input image is processed in three geometric variants: the original, a horizontal flip, and a vertical flip. Each variant is independently preprocessed through the standard transform pipeline and passed through the model. The classification logits from all three variants are averaged element-wise before softmax, producing more robust and stable predictions by reducing sensitivity to minor image orientation and spatial artifacts. The segmentation mask is taken from the original (non-augmented) input only.

\textbf{Out-of-Distribution (OOD) Detection:} After computing the softmax probability distribution, Shannon Entropy is calculated as $H = -\sum_{i} p_i \log p_i$ and normalized by $\log(K)$ where $K{=}2$ is the number of classes. This produces a value between 0 (maximum confidence) and 1 (maximum uncertainty). If the normalized entropy exceeds a threshold of 0.85, the image is flagged as Out-of-Distribution and the system returns a rejection response with a warning message, preventing unreliable predictions on non-fundus images such as photographs of cats, buildings, or corrupted medical scans.

\textbf{Temperature Scaling for Calibration:} Modern deep networks commonly produce overconfident predictions where softmax probabilities saturate very close to 0.0 or 1.0 even for ambiguous inputs. To address this known calibration issue, we apply a post-hoc temperature scaling factor of $T{=}2.5$ to the raw logits before softmax computation during inference: $p_i = \text{softmax}(z_i / T)$. This softens the probability distribution to produce confidence scores that more faithfully reflect the model's true predictive uncertainty, producing clinically realistic probability estimates rather than misleading 100\% confidence values.

\textbf{Human-in-the-Loop Diagnostic Threshold:} The binary classification decision boundary is exposed as a user-adjustable parameter through the frontend slider (default $= 0.5$, range $= 0.1$ to $0.9$, step $= 0.05$). A lower threshold increases sensitivity (detecting more true Glaucoma cases at the cost of more false positives), making the system suitable for population-level screening where missing a case is unacceptable. A higher threshold increases specificity, reducing false alarms and making the system suitable for confirmatory diagnosis. This allows clinicians to dynamically adapt the system's operating point without retraining.

\subsection{Explainable AI (XAI) Methods}

\textbf{Grad-CAM~\cite{selvaraju2017}:} We apply Gradient-weighted Class Activation Mapping to the CBAM module at the bottleneck of the active inference model. The gradients of the predicted class score with respect to the CBAM output feature maps are globally average-pooled to obtain importance weights for each channel. The weighted combination of feature maps produces a coarse localization heatmap that is upsampled to input resolution using bilinear interpolation and overlaid on the original fundus image with user-adjustable opacity (default 60\%). This allows clinicians to visually verify that the model's diagnostic attention is concentrated on the optic nerve head region rather than irrelevant background areas.

\textbf{Integrated Gradients~\cite{sundararajan2017}:} As a complementary XAI method providing a fundamentally different perspective, we compute Integrated Gradients using Meta's Captum library with 50 interpolation steps. Starting from a black baseline image (representing a neutral input with no diagnostic information), the method numerically approximates the path integral of gradients along the straight line from baseline to input: $\text{IG}_i(x) = (x_i - x'_i) \times \int_0^1 \frac{\partial F(x' + \alpha(x-x'))}{\partial x_i} d\alpha$. The absolute values of the resulting per-pixel attributions are averaged across the three RGB channels and normalized to the $[0, 1]$ range, producing a fine-grained importance map where each bright pixel represents a specific location that actively contributed to the model's prediction. The final visualization is a 50\% blend of the attribution heatmap and the original image.

% ============================================================
% V. EXPERIMENTAL SETUP
% ============================================================
\section{Experimental Setup}

\subsection{Dataset}
The ACRIMA dataset~\cite{diazpinto2019} contains 705 color retinal fundus images, each independently labeled by experienced Glaucoma specialists as either Normal or Glaucomatous. All images are pre-cropped to center on the optic disc region. Glaucomatous images are identified by the presence of ``\_g\_'' in their filenames. The dataset distribution and split are detailed in Table~\ref{tab:dataset}.

\begin{table}[h]
\centering
\caption{ACRIMA Dataset Distribution and Split}
\label{tab:dataset}
\begin{tabular}{lccc}
\toprule
\textbf{Split} & \textbf{Normal} & \textbf{Glaucoma} & \textbf{Total} \\
\midrule
Training (70\%)       & 216              & 277                & 493            \\
Validation (15\%)     & 46               & 59                 & 105            \\
Test (15\%)           & 47               & 60                 & 107            \\
\midrule
\textbf{Total}        & \textbf{309}     & \textbf{396}       & \textbf{705}   \\
\bottomrule
\end{tabular}
\end{table}

\subsection{Libraries and Tools}
The complete technology stack includes: Python~3.10 as the programming language; PyTorch~2.x and torchvision for deep learning model construction, training, and pre-trained weight loading; OpenCV (headless variant) for image processing, contour detection, and segmentation mask overlay generation; the pytorch-grad-cam library for Grad-CAM computation; Meta's Captum library for Integrated Gradients attribution; SciPy for scientific computing utilities; FastAPI with Uvicorn as the asynchronous ASGI backend server; React~19 with Vite~8 for the frontend single-page application; and scikit-learn for evaluation metric computation.

\subsection{Hardware and Software Environment}
All model training (both Teacher and Student stages) was performed on Google Colab using an NVIDIA Tesla T4 GPU with 16\,GB of dedicated VRAM and CUDA acceleration. The total combined training time for both the Teacher (15 epochs) and Student distillation (15 epochs) was approximately 90 minutes. Local inference, backend API serving, and frontend development were performed on a machine with an Intel Core i9 13th Generation processor and 16\,GB of DDR5 system RAM running Windows~11.

\subsection{Evaluation Metrics}
Classification performance was evaluated using five standard metrics computed on the held-out test set that was never exposed during training or model selection: Accuracy (overall correct prediction rate), Sensitivity/Recall (true positive rate --- ability to detect Glaucoma when present), Specificity (true negative rate --- ability to correctly identify healthy eyes), F1-Score (harmonic mean of Precision and Recall), and Area Under the ROC Curve (AUC-ROC, measuring discriminative performance across all thresholds).

% ============================================================
% VI. RESULTS
% ============================================================
\section{Final Results and Comprehensive Comparison}

\subsection{Performance Across All Project Phases}
Table~\ref{tab:comparison} presents a comprehensive comparison of the system's evolution across all four development phases, encompassing both quantitative metrics and qualitative capability improvements.

\begin{table}[h]
\centering
\caption{Comprehensive Comparison Across All Project Phases}
\label{tab:comparison}
\renewcommand{\arraystretch}{1.2}
\begin{tabular}{@{}lcccc@{}}
\toprule
\textbf{Metric} & \textbf{P1} & \textbf{P2} & \textbf{P3} & \textbf{P4 (Final)} \\
\midrule
Architecture     & Custom CNN       & EffNet-B0       & EffNet-B0       & T-S + CBAM        \\
Accuracy         & 72.5\%           & 99.05\%         & 99.05\%         & 95.5\%            \\
F1-Score         & 0.68             & 0.98            & 0.99            & 0.95              \\
Model Size       & 15\,MB           & 20\,MB          & 20\,MB          & 24\,MB            \\
Explainability   & None             & None            & Grad-CAM        & Grad-CAM + IG     \\
Segmentation     & No               & No              & No              & Yes (Cup/Disc)    \\
OOD Detection    & No               & No              & No              & Yes (Entropy)     \\
TTA              & No               & No              & No              & Yes (3-variant)   \\
Threshold Adj.   & No               & No              & No              & Yes (0.1--0.9)    \\
Deployment       & None             & None            & Streamlit       & FastAPI + React   \\
\bottomrule
\end{tabular}
\end{table}

\subsection{Critical Analysis of Results}
The apparent numerical accuracy decrease from Phase~2/3 (99.05\%) to Phase~4 (95.5\%) warrants careful and nuanced interpretation rather than surface-level comparison. The Phase~2/3 model was a simple binary classifier with a frozen EfficientNet-B0 backbone fine-tuned exclusively on the single ACRIMA dataset, achieving near-perfect but potentially overfitted performance. With only 705 images and a frozen pre-trained backbone, the model effectively learned to match the specific visual patterns of this particular dataset's imaging conditions and patient demographics.

The Phase~4 system, by contrast, is a fundamentally more complex multi-task architecture where the classification objective competes for shared representational capacity with the segmentation objective. The segmentation task acts as an implicit regularizer, preventing the classification head from exploiting dataset-specific shortcuts. Furthermore, the Student model is not trained through direct supervision on the original data but rather through knowledge distillation from the Teacher, introducing an additional knowledge compression step that inherently trades marginal task-specific accuracy for improved generalization capability.

The Phase~4 system compensates for this marginal accuracy difference through substantial qualitative improvements that are critical for real-world clinical deployment: dual XAI visualizations that satisfy clinical transparency and trust requirements, OOD detection that prevents dangerous unreliable predictions on out-of-domain inputs, TTA that stabilizes predictions against capture-time orientation artifacts, adjustable diagnostic thresholds that give clinicians control over the sensitivity-specificity operating point, and a professional web-based deployment that can serve real-time predictions in actual clinical workflows.

% ============================================================
% VII. MODEL OPTIMIZATION
% ============================================================
\section{Model Optimization and Justification}

\subsection{Why the Final Model Was Selected}
The selection of EfficientNet-B0 as the deployed Student model was driven by the requirement for an optimal balance between three competing objectives: model accuracy, inference speed, and deployment footprint. Among the architectures evaluated during the project lifecycle, VGG-16 (138M parameters, approximately 550\,MB) was immediately rejected due to its excessive size making web deployment completely impractical. ResNet-50 (25.6M parameters, 153\,MB weight file) demonstrated strong feature extraction capabilities and was therefore retained as the Teacher model, but its large file size and slower inference speed make it unsuitable as the primary deployed model for resource-constrained environments. EfficientNet-B0, with only 5.3M backbone parameters and a total deployed size of 24\,MB (including the CBAM module and segmentation decoder), is approximately $6{\times}$ smaller than the Teacher while retaining the distilled diagnostic knowledge through the Knowledge Distillation process.

\subsection{Impact of Hyperparameter Tuning}
Several hyperparameters were carefully tuned to optimize the distillation process. The distillation temperature ($T{=}4.0$) controls how much the Teacher's probability distribution is softened before being presented to the Student. A temperature of 4.0 was selected to sufficiently soften the distribution to expose meaningful inter-class probability relationships (the ``dark knowledge'') without excessively flattening it to the point where class discrimination information is lost. The inference-time temperature scaling ($T{=}2.5$) was calibrated to produce confidence scores that avoid the clinically misleading extreme 0\%/100\% probability saturation that is common in deep networks, resulting in realistic probability estimates that a clinician can meaningfully interpret. The Joint Loss weighting during distillation ($\alpha{=}0.5$ for classification, $\beta{=}1.0$ for segmentation) was designed to prioritize segmentation learning, as the segmentation task provides stronger and more informative spatial supervision signals that force the Student to develop richer internal feature representations of the optic disc and cup anatomy.

\subsection{Impact of CBAM Attention}
The CBAM module positioned at the bottleneck serves as a learned feature filter that actively suppresses diagnostically irrelevant background noise (retinal vasculature patterns, imaging artifacts, peripheral retinal features) while amplifying features related to the optic nerve head structure. By applying channel attention first (identifying which of the 2048/1280 feature maps encode relevant diagnostic information) followed by spatial attention (producing a pixel-level mask highlighting where the diagnostic features are located within the spatial extent of the feature map), CBAM ensures that both the downstream classification head and the Grad-CAM visualization focus their analysis on the correct anatomical region. This attention-guided focusing is analogous to the clinical workflow of an ophthalmologist who naturally directs their diagnostic gaze to the optic disc area when assessing a fundus photograph for Glaucoma.

% ============================================================
% VIII. NOVELTY
% ============================================================
\section{Novelty and Final Contribution}

\subsection{Innovative Aspects}
VisionGuard AI introduces a unique integration of six orthogonal technical contributions that, individually, have been explored in the literature but have not been previously combined into a single Glaucoma screening pipeline:

\begin{enumerate}
    \item \textbf{Knowledge Distillation for Model Compression:} Compressing a heavy ResNet-50 Teacher into a lightweight EfficientNet-B0 Student enables deployment on resource-constrained clinical hardware without sacrificing the Teacher's learned diagnostic representations.
    \item \textbf{Multi-Task Learning with Joint Objectives:} Simultaneous classification and segmentation forces the network to learn spatially-aware features that encode both diagnostic class boundaries and anatomical structure, providing implicit regularization.
    \item \textbf{CBAM Bottleneck Attention:} The attention module serves as a learned diagnostic focus mechanism, ensuring the network concentrates on the optic nerve head --- the clinically relevant anatomical landmark for Glaucoma assessment.
    \item \textbf{Dual Explainable AI Outputs:} Combining Grad-CAM (coarse regional heatmap) with Integrated Gradients (fine-grained pixel attributions) provides clinicians with two complementary and independently verifiable perspectives on the model's diagnostic reasoning.
    \item \textbf{Clinical Safety Framework:} OOD detection via entropy thresholding, adjustable diagnostic thresholds for sensitivity-specificity tuning, TTA for prediction stabilization, and temperature-calibrated confidence scores collectively create a safety layer suitable for clinical deployment.
    \item \textbf{Full-Stack Production Deployment:} A complete FastAPI REST backend with a React-based glassmorphism clinical dashboard, featuring interactive diagnostic controls (model selection, threshold slider, heatmap opacity, TTA/XAI toggles), demonstrates deployment readiness far beyond a research prototype.
\end{enumerate}

\subsection{Real-World Applicability}
The final Student model's compact footprint (24\,MB) and fast inference time (sub-200\,ms on CPU) make it suitable for deployment in mobile health (mHealth) clinics in rural and semi-urban areas where internet bandwidth is limited, dedicated GPU hardware is unavailable, and computing devices are basic. The RESTful API architecture with stateless request handling allows the system to be integrated with minimal engineering effort into existing hospital information systems, telemedicine platforms, electronic health record (EHR) systems, or standalone mobile applications. The Human-in-the-Loop threshold adjustment feature ensures that clinicians retain ultimate diagnostic authority, with the AI system serving as a decision support tool rather than a replacement for clinical judgment.

% ============================================================
% IX. DEPLOYMENT
% ============================================================
\section{Deployment Considerations}

\subsection{System Architecture}
The deployed system follows a decoupled client-server architecture designed for flexibility and scalability.

\textbf{Backend (FastAPI + Uvicorn):} The FastAPI server exposes three endpoints: (1)~\texttt{POST /api/screen} --- the primary inference endpoint accepting multipart file uploads with configurable query parameters for model selection (\texttt{teacher}/\texttt{student}), diagnostic threshold (0.0--1.0), TTA toggle, and XAI comparison toggle; (2)~\texttt{GET /api/metrics} --- returns benchmark performance metrics; (3)~\texttt{GET /api/health} --- returns system health status and available feature list. Both Teacher and Student model weights are loaded into memory at server startup, enabling instant model switching without reload latency. The server returns a comprehensive JSON response containing: diagnosis label, Glaucoma probability, calibrated confidence score, entropy score, Cup-to-Disc Ratio, medical recommendation text, and base64-encoded images for Grad-CAM heatmap, Integrated Gradients attribution map, and segmentation overlay.

\textbf{Frontend (React + Vite):} The React single-page application provides a premium glassmorphism-styled clinical dashboard with: (a)~drag-and-drop image upload with preview, (b)~model selection dropdown (EfficientNet-B0 Student or ResNet-50 Teacher), (c)~diagnostic threshold slider with real-time re-analysis on change, (d)~heatmap overlay intensity slider, (e)~TTA and XAI comparison toggles, (f)~four color-coded diagnostic metric cards (Diagnosis, Confidence, CDR, Entropy), (g)~an information row displaying model type, TTA status, threshold, and Glaucoma probability, (h)~a side-by-side XAI comparison viewer with adjustable-opacity Grad-CAM and Integrated Gradients overlays, and (i)~a color-coded medical recommendation banner.

\subsection{Input-Output Flow}
The complete diagnostic flow proceeds as follows:
\begin{enumerate}
    \item The clinician uploads a retinal fundus image via drag-and-drop or file browser.
    \item The React frontend displays a preview and sends the image with all selected parameters to the FastAPI backend via an HTTP POST request.
    \item The backend preprocesses the image through the transform pipeline, runs inference through the selected model (optionally with TTA), applies temperature scaling, computes the softmax probability distribution, evaluates OOD entropy, generates Grad-CAM and Integrated Gradients heatmaps, processes segmentation masks and computes CDR, and determines the medical recommendation.
    \item The backend returns a comprehensive JSON response containing all computed results and base64-encoded visualization images.
    \item The React frontend parses the response and renders interactive diagnostic results with color-coded metrics, XAI comparison viewer, and clinical recommendation.
\end{enumerate}

\subsection{Scalability and Integration}
FastAPI is built on the Starlette ASGI framework and natively supports asynchronous request handling, making it capable of serving concurrent multi-user requests efficiently. The stateless API design, where no session information is maintained between requests, allows straightforward horizontal scaling behind a load balancer or reverse proxy. The lightweight Student model's 24\,MB memory footprint enables multiple model instances to be loaded simultaneously on a single server, and the base64-encoded image response format ensures compatibility with any HTTP client regardless of platform.

% ============================================================
% X. LIMITATIONS
% ============================================================
\section{Limitations and Challenges}

\subsection{Data Limitations}
The ACRIMA dataset contains only 705 images sourced from a single clinical institution and patient demographic. This limited sample size and lack of diversity across ethnicities, imaging hardware, and clinical environments may result in reduced generalization when the system encounters retinal images with visual characteristics significantly different from the training distribution. Furthermore, the ACRIMA dataset does not include ground-truth segmentation masks for the optic disc and cup boundaries. Consequently, the segmentation decoder was trained with approximate dummy masks derived from pixel intensity thresholds rather than expert-drawn annotations, meaning the CDR computation is not yet fully supervised by clinical ground truth.

\subsection{Remaining Technical Weaknesses}
The model may produce unreliable predictions on fundus images captured under challenging clinical conditions, including eyes with severe cataracts that obstruct the view of the optic disc, images with significant flash glare or reflection artifacts, and photographs captured with consumer smartphone cameras rather than dedicated clinical fundus cameras. These conditions introduce visual artifacts and imaging characteristics that are entirely absent from the clean, standardized ACRIMA training data. Additionally, the system performs only binary classification (Normal vs.\ Glaucoma) and does not provide clinically actionable severity grading (e.g., Early, Moderate, or Advanced Glaucoma), which is necessary for appropriate treatment planning and patient management.

\subsection{Computational Constraints}
Training the complete Teacher-Student Knowledge Distillation pipeline required careful GPU memory management on the 16\,GB Colab T4 GPU. During the Student distillation stage, the simultaneous forward pass through both the frozen Teacher and the trainable Student, combined with the computation of four distinct loss components (ground-truth classification, ground-truth segmentation, classification KD, segmentation KD), placed significant demands on GPU memory. This necessitated a reduced batch size of 16 compared to larger batch sizes typically used in deep learning research.

% ============================================================
% XI. FUTURE WORK
% ============================================================
\section{Future Work}

Several concrete avenues for future development have been identified to extend VisionGuard AI's capabilities and clinical readiness:

\begin{enumerate}
    \item \textbf{Multi-Dataset Validation:} Training and evaluating the system on additional benchmark datasets including REFUGE~\cite{orlando2020}, DRISHTI-GS, and ORIGA to validate cross-dataset generalization and identify potential institutional or demographic biases.
    
    \item \textbf{Fully Supervised Segmentation:} Acquiring datasets with expert-drawn ground-truth optic cup and disc masks (such as DRISHTI-GS, which provides pixel-level annotations) to replace the current approximate masks and enable clinically accurate, fully supervised CDR computation.
    
    \item \textbf{Large-Scale Training:} Leveraging the AIROGS dataset, which contains over 113,000 retinal images from a diverse patient population, to train a significantly more robust and generalizable model suitable for production clinical deployment at scale.
    
    \item \textbf{Mobile Application Deployment:} Converting the trained Student model to optimized formats such as ONNX or TorchScript for deployment on Android and iOS devices, enabling point-of-care Glaucoma screening using smartphone-attached portable fundus cameras in field settings.
    
    \item \textbf{Multi-Class Severity Grading:} Extending the classification head from binary (Normal/Glaucoma) to multi-class (Normal/Mild/Moderate/Severe) to provide clinically actionable severity assessments that can directly inform treatment decisions and referral urgency.
    
    \item \textbf{Multi-Modal Fusion:} Integrating Optical Coherence Tomography~(OCT) volumetric scan data alongside 2D fundus images in a dual-stream architecture, as recent literature consistently demonstrates that combined multi-modal imaging analysis produces the most reliable and comprehensive clinical assessment of Glaucoma progression.
\end{enumerate}

% ============================================================
% XII. CONCLUSION
% ============================================================
\section{Conclusion}

This paper presented VisionGuard AI, a complete research-to-deployment deep learning pipeline for automated Glaucoma screening from retinal fundus images. Over four iterative development phases spanning an entire academic semester, the system evolved from a baseline four-layer CNN achieving 72.5\% accuracy into a sophisticated Teacher-Student Knowledge Distillation framework incorporating CBAM attention mechanisms, multi-task learning with joint classification and segmentation objectives, dual Explainable AI visualizations via Grad-CAM and Integrated Gradients, and a comprehensive clinical safety framework including OOD detection, TTA, adjustable diagnostic thresholds, and temperature-calibrated confidence scores. The final deployed Student model (EfficientNet-B0 with CBAM, 24\,MB) achieves competitive classification performance while providing transparent, visually verifiable predictions that directly address the ``black-box'' trust barrier that remains the primary obstacle to clinical adoption of AI-assisted diagnostics. The system's lightweight architecture, sub-200\,ms CPU inference latency, and full-stack web deployment through FastAPI and React demonstrate that research-grade medical AI can be made practical, accessible, and clinically trustworthy. VisionGuard AI represents a meaningful contribution toward bridging the gap between academic deep learning research and deployable, real-world healthcare solutions for underserved populations where access to specialist ophthalmological care remains critically limited.

% ============================================================
% REFERENCES
% ============================================================
\balance
\begin{thebibliography}{10}

\bibitem{who2019}
World Health Organization, ``World Report on Vision,'' WHO, Geneva, Switzerland, 2019.

\bibitem{li2018}
Z.~Li, Y.~He, S.~Keel, W.~Meng, R.~Chang, and M.~He, ``Efficacy of a Deep Learning System for Detecting Glaucomatous Optic Neuropathy Based on Color Fundus Photographs,'' \textit{Ophthalmology}, vol.~125, no.~8, pp.~1199--1206, 2018.

\bibitem{diazpinto2019}
A.~Diaz-Pinto, S.~Morales, V.~Naranjo, T.~K{\"o}hler, J.~M.~Mossi, and A.~Navea, ``CNNs for Automatic Glaucoma Assessment Using Fundus Images: An Extensive Validation,'' \textit{BioMedical Engineering OnLine}, vol.~18, no.~29, 2019.

\bibitem{orlando2020}
J.~I.~Orlando, H.~Fu, J.~B.~Breda, K.~van Keer, D.~R.~Bathula, A.~Diaz-Pinto, \textit{et~al.}, ``REFUGE Challenge: A Unified Framework for Evaluating Automated Methods for Glaucoma Assessment from Fundus Photographs,'' \textit{Medical Image Analysis}, vol.~59, 2020.

\bibitem{tan2019}
M.~Tan and Q.~V.~Le, ``EfficientNet: Rethinking Model Scaling for Convolutional Neural Networks,'' in \textit{Proc.\ Int.\ Conf.\ Machine Learning (ICML)}, 2019.

\bibitem{he2016}
K.~He, X.~Zhang, S.~Ren, and J.~Sun, ``Deep Residual Learning for Image Recognition,'' in \textit{Proc.\ IEEE Conf.\ Computer Vision and Pattern Recognition (CVPR)}, pp.~770--778, 2016.

\bibitem{selvaraju2017}
R.~R.~Selvaraju, M.~Cogswell, A.~Das, R.~Vedantam, D.~Parikh, and D.~Batra, ``Grad-CAM: Visual Explanations from Deep Networks via Gradient-Based Localization,'' in \textit{Proc.\ IEEE Int.\ Conf.\ Computer Vision (ICCV)}, pp.~618--626, 2017.

\bibitem{sundararajan2017}
M.~Sundararajan, A.~Taly, and Q.~Yan, ``Axiomatic Attribution for Deep Networks,'' in \textit{Proc.\ Int.\ Conf.\ Machine Learning (ICML)}, pp.~3319--3328, 2017.

\bibitem{woo2018}
S.~Woo, J.~Park, J.~Lee, and I.~S.~Kweon, ``CBAM: Convolutional Block Attention Module,'' in \textit{Proc.\ European Conf.\ Computer Vision (ECCV)}, pp.~3--19, 2018.

\bibitem{hinton2015}
G.~Hinton, O.~Vinyals, and J.~Dean, ``Distilling the Knowledge in a Neural Network,'' in \textit{NIPS Deep Learning and Representation Learning Workshop}, 2015.

\end{thebibliography}

\end{document}
